\paragraph{Rewriting Rules}
Given a logical plan as above, there are multiple ways in which we can construct a physical plan.
We can use different join orders, decide when selection happens and specific implementations of operators.

The following rules were discussed in the lectures:
\begin{enumerate}[label=\textbf{R\arabic*}]
    \item Conjunctive selection operators can be deconstructed into a sequence of individual selections:
          \[
              \sigma_{\theta_1 \land \theta_2} = \sigma_{\theta_1}(\sigma_{\theta_2}(E))
          \]
          This is useful to split selections to push them earlier in the execution plan (reduces the number of elements to process subsequently)
    \item Selection operators are commutative:
          \[
              \sigma_{\theta_1}(\sigma_{\theta_2}(E)) = \sigma_{\theta_2}(\sigma_{\theta_1}(E))
          \]
          This is useful to allow to first put the one that e.g. has an index.
    \item Only the last in a sequence of projections is needed, others can be omitted
          \[
              \Pi_{t_1}(\Pi_{t_2}(E)) = \Pi_{t_1}(E)
          \]
\end{enumerate}
